\subsection{Byte-Oriented Data Streams}

Text files pass through an encoding and decoding boundary. Binary files do not. In binary mode, Python treats the file as a stream of bytes. The program does not ask Python to interpret those bytes as characters, lines, words, or encoded text.

This is the correct model for images, compressed archives, executable files, audio data, network packets, and any format where each byte value matters directly. A binary file is not ``text with strange characters''. It is byte data.

\subsubsection{Binary Mode as Raw Byte Transfer Without Text Decoding}

Binary mode is selected by adding \texttt{b} to the mode string. The most common binary modes are \texttt{"rb"} for reading bytes and \texttt{"wb"} for writing bytes.

In binary mode, \texttt{write()} expects a bytes-like object, not a \texttt{str}. Likewise, \texttt{read()} returns \texttt{bytes}, not \texttt{str}.

\begin{verbatim}
path = "binary_mode_demo.bin"

data = b"Binary mode does not decode text. 42.\n"

print("Byte data before writing:")
print(f"data       = {data}")
print(f"type(data) = {type(data)}")
print(f"len(data)  = {len(data)}")

print()
print("Selected names from dir(data):")
for name in ["decode", "hex", "split", "replace", "startswith", "count"]:
    print(f"{name:12} -> {name in dir(data)}")

with open(path, "wb") as file_obj:
    print()
    print("Binary file object while writing:")
    print(f"file_obj       = {file_obj}")
    print(f"type(file_obj) = {type(file_obj)}")
    print(f"file_obj.mode  = {file_obj.mode}")

    print()
    print("Selected names from dir(file_obj):")
    for name in ["write", "read", "close", "flush", "seek", "tell", "mode"]:
        print(f"{name:10} -> {name in dir(file_obj)}")

    written_count = file_obj.write(data)

print()
print(f"write() returned: {written_count}")
print("In binary mode, this is the number of bytes written.")

with open(path, "rb") as file_obj:
    restored_data = file_obj.read()

print()
print("Byte data read back:")
print(f"restored_data       = {restored_data}")
print(f"type(restored_data) = {type(restored_data)}")
print(f"data == restored_data -> {data == restored_data}")
\end{verbatim}

The object written to the file is a \texttt{bytes} object. The object read back from the file is also a \texttt{bytes} object. Python does not decode the file content into text because the file was opened in binary mode.

Passing a string directly to a binary file is an error.

\begin{verbatim}
path = "binary_type_error_demo.bin"

text = "D'oh! This is a str, not bytes."

with open(path, "wb") as file_obj:
    try:
        file_obj.write(text)
    except TypeError as err:
        print("Writing str to a binary file failed.")
        print(f"type(err) = {type(err)}")
        print(f"err       = {err}")

    fixed_data = text.encode("utf-8")
    file_obj.write(fixed_data)

print()
print("Corrected by explicit encoding:")
print(f"fixed_data       = {fixed_data}")
print(f"type(fixed_data) = {type(fixed_data)}")
\end{verbatim}

The failure is useful. It prevents the program from silently choosing an encoding. If the programmer wants text to become bytes, the conversion should be explicit.

\subsubsection{The \texttt{bytes} Object: Immutable Byte Sequences Distinct from \texttt{str}}

A \texttt{bytes} object is an immutable sequence of byte values. Each element is an integer from \texttt{0} to \texttt{255}. This is different from a \texttt{str}, whose elements are Unicode characters.

\begin{verbatim}
text = "Spam and eggs."
data = b"Spam and eggs."

print("Text object:")
print(f"text       = {text}")
print(f"type(text) = {type(text)}")
print(f"len(text)  = {len(text)}")

print()
print("Bytes object:")
print(f"data       = {data}")
print(f"type(data) = {type(data)}")
print(f"len(data)  = {len(data)}")

print()
print("Indexing:")
first_char = text[0]
first_byte = data[0]

print(f"text[0]       = {first_char!r}")
print(f"type(text[0]) = {type(first_char)}")

print(f"data[0]       = {first_byte}")
print(f"type(data[0]) = {type(first_byte)}")

print()
print("Slicing:")
print(f"text[0:4]       = {text[0:4]!r}")
print(f"type(text[0:4]) = {type(text[0:4])}")

print(f"data[0:4]       = {data[0:4]}")
print(f"type(data[0:4]) = {type(data[0:4])}")
\end{verbatim}

The difference is precise. Indexing a string returns a one-character \texttt{str}. Indexing a \texttt{bytes} object returns an \texttt{int}. Slicing a \texttt{bytes} object returns another \texttt{bytes} object.

A \texttt{bytes} object can be inspected and transformed without becoming text.

\begin{verbatim}
data = b"Nobody expects the Spanish Inquisition!"

print("Original bytes:")
print(data)

print()
print("Selected bytes methods:")
for name in ["hex", "decode", "replace", "split", "startswith", "endswith"]:
    print(f"{name:12} -> {name in dir(data)}")

print()
print("Byte-level operations:")
print(f"data.startswith(b'Nobody') -> {data.startswith(b'Nobody')}")
print(f"data.endswith(b'!')        -> {data.endswith(b'!')}")
print(f"data.count(b'e')           -> {data.count(b'e')}")
print(f"data.replace(b'Nobody', b'Everybody') -> {data.replace(b'Nobody', b'Everybody')}")

print()
hex_text = data.hex()
print("Hexadecimal view:")
print(f"hex_text       = {hex_text}")
print(f"type(hex_text) = {type(hex_text)}")
\end{verbatim}

The \texttt{hex()} method returns a textual hexadecimal representation of the byte sequence. This is useful for inspection because arbitrary bytes are not always printable as meaningful text.

The \texttt{decode()} method crosses from bytes back to text, but only if the bytes represent valid text under the chosen encoding.

\begin{verbatim}
data = b"Serenity now!"

text = data.decode("utf-8")

print("Decoded bytes:")
print(f"text       = {text}")
print(f"type(text) = {type(text)}")
\end{verbatim}

The practical rule is: \texttt{bytes} is not a rough synonym for \texttt{str}. It is a separate object type for byte sequences. Convert explicitly when crossing between text and byte data.

\subsubsection{The \texttt{bytearray} Object: Mutable Byte Buffers for Incremental Modification}

A \texttt{bytes} object is immutable. Once created, its byte values cannot be changed in place. A \texttt{bytearray} is the mutable counterpart. It stores byte values like \texttt{bytes}, but allows modification.

\begin{verbatim}
data = b"Chuck Norris fact: 41."

print("Immutable bytes object:")
print(f"data       = {data}")
print(f"type(data) = {type(data)}")

print()
try:
    data[-2] = 50
except TypeError as err:
    print("Changing bytes in place failed.")
    print(f"type(err) = {type(err)}")
    print(f"err       = {err}")
\end{verbatim}

To modify byte data incrementally, create a \texttt{bytearray}.

\begin{verbatim}
buf = bytearray(b"Chuck Norris fact: 41.")

print("Mutable bytearray object:")
print(f"buf       = {buf}")
print(f"type(buf) = {type(buf)}")
print(f"len(buf)  = {len(buf)}")

print()
print("Selected names from dir(buf):")
for name in ["append", "extend", "insert", "pop", "replace", "decode", "hex"]:
    print(f"{name:10} -> {name in dir(buf)}")

print()
print("Before modification:")
print(buf)

buf[-2] = ord("2")

print("After modification:")
print(buf)

text = buf.decode("utf-8")

print()
print("Decoded modified buffer:")
print(f"text       = {text}")
print(f"type(text) = {type(text)}")
\end{verbatim}

The expression \texttt{ord("2")} returns the integer code value for the character \texttt{"2"}. A \texttt{bytearray} position stores an integer byte value, not a one-character string.

A \texttt{bytearray} can also grow.

\begin{verbatim}
buf = bytearray()

print("Empty bytearray:")
print(f"buf       = {buf}")
print(f"type(buf) = {type(buf)}")

buf.extend(b"D'oh")
buf.append(33)
buf.extend(b" 42")

print()
print("After incremental changes:")
print(f"buf = {buf}")

print()
print("Element inspection:")
for index, value in enumerate(buf):
    print(f"index {index:02d} -> value {value:3d}")

print()
print("As decoded text:")
print(buf.decode("utf-8"))
\end{verbatim}

The call \texttt{append(33)} adds one byte. The value \texttt{33} is the byte value for the exclamation mark in ASCII-compatible encodings. The call \texttt{extend(...)} adds several bytes.

For file output, both \texttt{bytes} and \texttt{bytearray} can be written in binary mode.

\begin{verbatim}
path = "bytearray_write_demo.bin"

buf = bytearray(b"The answer was ")
buf.extend(b"41")
buf[-1] = ord("2")
buf.extend(b".\n")

with open(path, "wb") as file_obj:
    file_obj.write(buf)

with open(path, "rb") as file_obj:
    restored_data = file_obj.read()

print("Restored binary data:")
print(f"restored_data       = {restored_data}")
print(f"type(restored_data) = {type(restored_data)}")

print()
print("Decoded for display:")
print(restored_data.decode("utf-8"))
\end{verbatim}

The \texttt{bytearray} is useful when byte data must be built or modified step by step. The \texttt{bytes} object is useful when the byte sequence should be fixed and safe from in-place mutation.

The practical distinction is direct:

\begin{itemize}
    \item \texttt{str}: immutable runtime text;
    \item \texttt{bytes}: immutable byte sequence;
    \item \texttt{bytearray}: mutable byte buffer.
\end{itemize}

Binary file mode works with byte-oriented objects. It does not perform text decoding or encoding unless the programmer explicitly asks for it.